\documentclass[10pt,a4paper]{article}
\usepackage[T1]{fontenc}
\usepackage[top=0.5in,bottom=0.5in,left=0.55in,right=0.55in]{geometry}
\usepackage{enumitem}
\usepackage[hidelinks]{hyperref}
\usepackage{titlesec}

\setlist[itemize]{leftmargin=1.2em,noitemsep,topsep=2pt,parsep=0pt,partopsep=0pt}
\titleformat{\section}{\normalsize\bfseries}{}{0em}{}[\titlerule]

\begin{document}
\pagestyle{empty}

%====================
% HEADER
%====================
\begin{center}
    {\Large \textbf{Learned Regularization for Sparse-View CT:}}\\
    {\large Accuracy, Lesion Detectability and Distribution Shift of a\\
    Plug-and-Play Prior against Tuned Total Variation}
\end{center}
%{\large \textbf{When does a learned prior help sparse-view CT, and when does it not?}}\\[0.2cm]
Mohammed Abo Oraig \quad\textbar\quad \href{mailto:mohammed-abo-oraig@etud.univ-tln.fr}{\underline{mohammed-abo-oraig@etud.univ-tln.fr}} \\
Code: \href{https://github.com/AboOraig/learned-regularization}{\underline{github.com/AboOraig/learned-regularization}}

\section*{Question}
Plug-and-Play (PnP) reconstruction with a learned denoiser is often reported to beat classical regularisers such as total variation (TV) in PSNR.
This study asks three questions under controlled conditions: \textbf{Q1} does the advantage survive a change in image distribution,
\textbf{Q2} can the method create lesion-like structure that the ground truth does not support, and \textbf{Q3} do the conclusions hold across noise levels and acquisition settings?

\section*{Approach}
A noise-conditioned DnCNN is used as the prior in PnP proximal gradient and compared with FBP, Tikhonov and TV (unconstrained and box-constrained), each tuned on validation data.
Three stages of increasing realism: 28$\times$28 digits with Gaussian noise; Poisson noise with a matrix-free operator validated against the dense one;
and LoDoPaB-CT slices (362$\times$362, patient-disjoint split, simulated low-dose measurements, 30/60/120 of 1000 angles, three doses, 40 test patients).
Beyond PSNR/SSIM the study uses (i) a singular-vector decomposition of the error into well- and weakly-determined components, (ii) controlled lesion insertion with a location-known observer (paired AUC, lesion retention, false blobs),
and (iii) controls against evaluation artefacts: the inverse crime, a box-constrained TV baseline, independent denoiser seeds and a matched-prior control.
Statistics are paired; bootstrap intervals over patients on CT.

\section*{Main findings}
\begin{itemize}
\item \textbf{The advantage is real but conditional.} On digits the prior is about 3.4\,dB above a box-constrained TV at moderate noise; it shrinks at both ends of the noise range, and near $\sigma=0$ the ranking depends on the data model and the TV solver.
\item \textbf{It reverses under far distribution shift, and the cause is prior mismatch.} An MNIST-trained prior is $3.10\pm0.30$\,dB (Fashion-MNIST) and $4.33\pm0.12$\,dB (phantom) \emph{below} unconstrained TV (mean $\pm$ s.d., three seeds; against box-constrained TV: $3.13$ and $4.55$\,dB), while a Fashion-trained prior is $2.21$\,dB \emph{above} TV on Fashion-MNIST with identical hyper-parameters.
\item \textbf{PSNR does not predict detectability.} In distribution the prior is $3$--$4$\,dB better in PSNR yet detects inserted lesions worse than TV because it attenuates them; under mismatch it creates more false lesions than box-constrained TV for seven of eight lesion types (TV's own false-blob count is unstable). Re-tuning trades one failure for the other.
\item \textbf{On CT slices the accuracy advantage disappears.} Against a TV run for 300 iterations, PnP differs by $-0.93$ to $+0.33$\,dB over nine dose$\times$angle conditions (SSIM lower in 8 of 9). Its lesion-detection AUC is higher than that TV's in all nine conditions (significant in six), but never significantly higher than plain FBP, and the PSNR-tuned TV is \emph{below} FBP at detection in all nine (seven significant).
\item \textbf{False structures on CT slices} mostly appear at the same locations in FBP, TV and PnP, which suggests that they come from the data rather than the prior; on non-anatomical phantoms PnP's PSNR is below TV in all six conditions.
\end{itemize}

\section*{What I contributed and what I did not claim}
Design and implementation of the pipeline (operators with exact adjoints, solvers, simulators, observer, unit tests), the experimental protocol and the analysis. The code is released with the stored results.
An AI assistant (Claude) helped develop the code and, later, to edit the text; the research questions, experiments, interpretation and conclusions are mine, as the paper's disclosure states.
The paper states its limits explicitly: one small denoiser and one CT training run, 28$\times$28 images for the thorough grid, 40 CT patients with one noise realisation per condition, synthetic lesions and no reader study, and no multiplicity correction.
Part of the CT analysis was developed after seeing results (documented in the paper); nothing was pre-registered.

\section*{Skills demonstrated}
Variational and learned regularisation for inverse problems; PnP and variational solvers (ADMM, primal--dual); PyTorch and matrix-free GPU implementation on a 4\,GB card; careful statistical evaluation and controls; reproducible research software.

\section*{Where this leads}
Open questions the study leaves: isolating why the detection sign differs between digits and CT; testing convergence-guaranteed or score-based priors under the same protocol; clinically low-contrast lesions with a reader study.
\end{document}